%% kapa: exam=1 round=36 year=2025 period=1 form=A subject=civil_law subject_ko=민법 seq=36 number=36
%% source: Q-Net 감정평가사 기출문제 (cst00309) artlSeq=5231525 file=raw/1008/5231525/2025년 제36회 감정평가사 1차시험 1교시 문제지 원본.pdf
%% conversion: kordoc 4.24.51 -> tools/gen1.py (auto); review=proofread-sonnet; flags=none
\begin{question}{36}
유치권에 관한 설명으로 옳지 않은 것은? (다툼이 있으면 판례에 따름)
\begin{choices}
  \item 유치권자는 유치물의 과실인 금전을 수취하여 다른 채권보다 먼저 그 채권의 변제에 충당할 수 있다.
  \item 유치권자가 소유자의 허락 없이 유치물을 임대한 경우 임차인은 소유자에게 임대차로 대항할 수 없다.
  \item 여러 필지의 토지에 유치권을 행사하는 자가 그 토지 중 일부에 대해 선관주의의무를 위반한 경우 모든 토지에 대한 유치권소멸청구가 인정된다.
  \item 유치권에 의한 경매로 유치물이 매각되는 경우 유치권자는 일반채권자와 동일한 순위로 배당을 받는다.
  \item 저당권이 설정된 건물에 대하여 경매개시결정 이전에 유치권이 성립한 때에는 유치권자는 경매절차의 매수인에게 대항할 수 있다.
\end{choices}
\end{question}
\answer{3}
